\section{Positive reference model and the true \texorpdfstring{one-$\chi$}{one-chi} defect}
\label{sec:model-floor}

The true logarithmic source is the pair
\[
H_{\log}(s)=-2P(s)+P(s-2/L),
\]
evaluated with the exact packet polarization and one-$\chi$ stationary phase.
We denote the resulting retained Hermitian operator by
\begin{equation}
\boxed{A_T^{\rm true}:=A_T[H_{\log};e(-ne^d)].}
\label{eq:true-logonly-operator}
\end{equation}
This notation emphasizes that the physical additive carrier $e(-ne^d)$ is
retained.  No claim is made that the untwisted scalar HLZ diagonal identifies
$A_T^{\rm true}$ with the comparison model below.  No critical-$J$ self
profile is included in either operator.

Let $C(x):=G_{\log}(1-x)$ on the retained interior lag interval.  If $q$ is
mean zero and
\[
Q(x):=\int_0^xq(t)\,dt,
\]
then two integrations by parts give the exact primitive identity
\begin{equation}
\boxed{
\langle K_Cq,q\rangle
=-2C'(0+)\|Q\|_2^2
-\iint C''(|x-y|)Q(x)\overline{Q(y)}\,dx\,dy.
}
\label{eq:primitive-identity}
\end{equation}
Define
\[
Q_\triangle(q):=2\|Q\|_2^2.
\]

On the central flat plateau translate the physical Fourier coordinate by its
fixed left endpoint $u_0$.  The flat packet synthesis is
\begin{equation}
\boxed{
(\mathcal U_T^{\rm flat}v)(u)
=L_{\rm pkt}^{-1/2}\sum_\nu
 e^{i\tau_\nu u_0}v_\nu e^{i\tau_\nu u}.
}
\label{eq:explicit-flat-analysis}
\end{equation}
Let $C_T$ be the stable-core restriction, put
\[
R_T^{\rm core}:=\frac{W_T}{L},
\qquad
\mathcal U_T^{\rm core}:=C_T\mathcal U_T^{\rm flat},
\]
and normalize by $r=u/L$:
\begin{equation}
\boxed{
(\widehat{\mathcal U}_T^{\rm core}v)(r)
:=L^{1/2}(\mathcal U_T^{\rm core}v)(Lr),
\qquad 0\le r\le R_T^{\rm core}.
}
\label{eq:normalized-core-analysis}
\end{equation}
The global reference form is the pullback, through this same core analysis,
of $Q_\triangle$.  Equivalently
\begin{equation}
\boxed{
\Gfr_T=(\mathcal W_T^{\rm glob})^*\mathcal W_T^{\rm glob},
}
\label{eq:global-reference-Gram}
\end{equation}
where $\mathcal W_T^{\rm glob}$ adjoins the primitive on
$0\le u\le W_T$ and the normalization $\sqrt2/L$.

Motivated by the candidate endpoint geometry in
Corollary~\ref{cor:frozen-logonly-lag-kernel}, define the positive reference
kernel
\[
K_{G_{\log}}(r,s)=G_{\log}(1-|r-s|),
\]
with
\begin{equation}
\boxed{
G_{\log}(\rho)
=-\frac76\rho^4+2\rho^3-\rho^2
+\frac{e^2}{4}\left[
(2\rho^2-4\rho+3)\sinh(2\rho)-\rho^2
\right].
}
\label{eq:Glog-closed}
\end{equation}
This is the log-only profile obtained from
\eqref{eq:Glog-from-source}.  In particular the hyperbolic term is
$\sinh(2\rho)$, not $\sinh^2\rho$.

For later certification we use the elementary rational enclosure
\begin{equation}
\sum_{j=0}^{m}\frac{x^j}{j!}
\le e^x\le
\sum_{j=0}^{m}\frac{x^j}{j!}
+\frac{x^{m+1}}{(m+1)!}\frac1{1-x/(m+2)}
\qquad(0\le x\le4,\ m+2>x).
\label{eq:exp-rational-bound}
\end{equation}

\begin{proposition}[Certified log-only spectral floor]
\label{prop:model-floor}
On every mean-zero stable core fibre,
\begin{equation}
\boxed{
K_{G_{\log}}\succeq0.98\,Q_\triangle.
}
\label{eq:model-floor}
\end{equation}
Consequently, for the reference packet pullback
\[
P_T:=A_T^{G_{\log}},
\]
one has
\begin{equation}
\boxed{P_T\succeq0.98\,\Gfr_T.}
\label{eq:logonly-model-floor}
\end{equation}
In the final perturbation argument we use only the much weaker fixed constant
\begin{equation}
\boxed{c_0:=0.02.}
\label{eq:PT-logonly-definition}
\end{equation}
This is a statement about the reference model $P_T$; it does not identify
$P_T$ with $A_T^{\rm true}$.
\end{proposition}

\begin{proof}
Extend a core input by zero to $[0,1]$.  Put
\[
g_{\log}(r):=G_{\log}''(r).
\]
Direct differentiation of \eqref{eq:Glog-closed} gives
\begin{equation}
\boxed{
\begin{aligned}
g_{\log}(r)={}&-14r^2+12r-2\\
&+e^2\left[(4r-4)\cosh(2r)
 +(2r^2-4r+4)\sinh(2r)-\frac12\right].
\end{aligned}}
\label{eq:glog-second}
\end{equation}
The cusp coefficient in \eqref{eq:primitive-identity} is
\begin{equation}
\boxed{
a_{\log}=2G_{\log}'(1)
=e^2(\cosh2-1)-\frac43>19.07.}
\label{eq:logonly-cusp}
\end{equation}

We first record the sign geometry of $g_{\log}$.  Write $x=1-r$.  A further
derivative gives
\[
g_{\log}'(r)
=-28r+12+2e^2\left[(2x^2+4)\cosh(2r)-6x\sinh(2r)\right].
\]
Since $\tanh(2r)\le2r=2(1-x)$,
\[
(2x^2+4)-6x\tanh(2r)
\ge14x^2-12x+4\ge\frac{10}{7}.
\]
Therefore
\begin{equation}
\boxed{
g_{\log}'(r)>-16+\frac{20}{7}e^2>4
\qquad(0\le r\le1),}
\label{eq:glog-monotone}
\end{equation}
so $g_{\log}$ is strictly increasing.  Applying
\eqref{eq:exp-rational-bound} with $m=64$ gives, by exact rational arithmetic,
\begin{equation}
 g_{\log}(0.5833)<-0.0039,
 \qquad
 g_{\log}(0.5834)>0.0015.
\label{eq:glog-root-certificate}
\end{equation}
Thus $g_{\log}$ has one and only one zero
\[
r_0\in(0.5833,0.5834).
\]

Let
\[
M_{\log}(x):=\int_0^1|g_{\log}(1-|x-y|)|\,dy.
\]
It is symmetric about $1/2$.  For $0\le x\le1/2$,
\begin{equation}
M_{\log}(x)
=\int_x^1|g_{\log}(r)|\,dr
 +\int_{1-x}^1|g_{\log}(r)|\,dr.
\label{eq:glog-row-sum}
\end{equation}
For $x<1-r_0$ its derivative is
\[
M_{\log}'(x)=g_{\log}(x)+g_{\log}(1-x).
\]
The same $m=64$ rational enclosure, applied on the indicated compact
intervals, gives
\begin{equation}
\begin{gathered}
\frac d{dx}\{g_{\log}(x)+g_{\log}(1-x)\}<-8
\qquad(0\le x\le0.417),\\
g_{\log}(0.1338)+g_{\log}(0.8662)>0.004,\\
g_{\log}(0.1339)+g_{\log}(0.8661)<-0.0015.
\end{gathered}
\label{eq:glog-row-max-certificate}
\end{equation}
Hence the unique maximum on this interval occurs at
$x_0\in(0.1338,0.1339)$.  For $1-r_0\le x\le1/2$, both arguments lie below
$r_0$, and strict increase of $g_{\log}$ gives
$M_{\log}'(x)=g_{\log}(x)-g_{\log}(1-x)<0$.  Thus $x_0$ is the global maximizer
up to the symmetry $x\mapsto1-x$.

Since $x_0<1-r_0$, integration of \eqref{eq:glog-row-sum} by the
antiderivative $G_{\log}'$ gives
\[
M_{\log}(x_0)
=G_{\log}'(x_0)+2G_{\log}'(1)
 -2G_{\log}'(r_0)-G_{\log}'(1-x_0).
\]
Applying the same rational Taylor enclosure on the certified intervals for
$r_0$ and $x_0$ yields the deliberately coarse bound
\begin{equation}
\boxed{\sup_{0\le x\le1}M_{\log}(x)<17.10.}
\label{eq:Glog-Schur}
\end{equation}
All inequalities in
\eqref{eq:glog-root-certificate}--\eqref{eq:Glog-Schur} are interval
consequences of \eqref{eq:exp-rational-bound}; no floating-point value enters
the proof.

Schur's test and \eqref{eq:primitive-identity} now give
\[
\langle K_{G_{\log}}q,q\rangle
>(19.07-17.10)\|Q\|_2^2
>1.97\|Q\|_2^2
>0.98Q_\triangle(q).
\]
The packet inequality follows by congruence through the stable-core analysis
map used in the definition of $P_T$.  This proves the proposition.
\end{proof}

\begin{hypothesis}[One-$\chi$ one-sided relative spectral defect]
\label{hyp:one-chi-relative-defect}
For the admissible retained packet data chosen in
Section~\ref{sec:multiplicity-endgame}, there is a deterministic sequence
$\lambda_T>0$ such that, with
\begin{equation}
\mathcal D_T
:=\Gfr_T^{-1/2}(A_T^{\rm true}-\lambda_TP_T)\Gfr_T^{-1/2},
\label{eq:true-model-relative-defect}
\end{equation}
one has
\begin{equation}
\boxed{
\beta_-(T;\lambda_T)
:=\frac{\tr[(\mathcal D_T)_-^2]}
{(0.98\lambda_T)^2N_T}
=O(L^{-2}).}
\label{eq:one-chi-relative-defect}
\end{equation}
Only the negative part is constrained.  In particular, no two-sided operator
norm or Hilbert--Schmidt approximation of $A_T^{\rm true}$ by
$\lambda_TP_T$ is assumed.
\end{hypothesis}

\begin{remark}[Unresolved input]
Hypothesis~\ref{hyp:one-chi-relative-defect} is not proved in this manuscript.
The surviving phase $e(-ne^d)$ is precisely why the scalar HLZ argument does
not establish it.  The normalization is chosen so that
$\beta_-(T;\lambda_T)=o(1)$ deletes only $o(N_T)$ spectral directions; the
displayed $O(L^{-2})$ rate gives the quantitative conditional conclusion in
Theorem~\ref{thm:main}.
\end{remark}

\begin{remark}[Numerical margin]
For orientation only, high-precision evaluation gives
\[
r_0=0.5833722208\ldots,
\qquad
\arg\max M_{\log}=0.13387266\ldots,
\qquad
\sup M_{\log}=17.07053067\ldots,
\]
and
\[
a_{\log}=19.07668558\ldots.
\]
The proof uses only the rational bounds above.
\end{remark}

\begin{lemma}[The whole-space floor survives restriction]
\label{lem:model-compression}
Let $V$ be any subspace of the finite packet coordinate space and let $P_V$ be
its orthogonal projection. Then
\[
P_VP_TP_V\succeq0.98\,P_V\Gfr_TP_V
\quad\text{on }V.
\]
\end{lemma}

\begin{proof}
For $v\in V$, $P_Vv=v$ and Proposition~\ref{prop:model-floor} applies directly.
\end{proof}

For a Hermitian lag form $B$ on the mean-zero fibre define
\begin{equation}
\|B\|_{\op,\triangle}
:=\sup_{\substack{q\ne0\\ \int q=0}}
\frac{|\langle Bq,q\rangle|}{Q_\triangle(q)}.
\label{eq:triangular-form-norm}
\end{equation}

\begin{proposition}[Frozen principal floor]
\label{prop:finiteT-principal-comparison}
On $V_{\rm mean}$,
\begin{equation}
\boxed{P_T\succeq0.98\,\Gfr_T\succeq c_0\Gfr_T,
\qquad c_0=0.02.}
\label{eq:principal-floor}
\end{equation}
The finite-height deformation is not absorbed into this positive operator; it
remains a separate relative-regular source.
\end{proposition}

\begin{proof}
This is Proposition~\ref{prop:model-floor} followed by restriction to
$V_{\rm mean}$.
\end{proof}
